% from ega3-1-fr.tex
\oldpage[0\textsubscript{III}]{38}
\begin{lemma}[11.5.2.4]
\label{0.11.5.2.4-fr}
Si $L^\bullet=(L^i)_{i\in\mathbb Z}$ est une résolution droite d'un objet $A$
de $\mathcal C$ telle que $\mathrm R^nT(L^i)=0$ pour tout $i\in\mathbb Z$ et
tout $n>0$, alors on a
$\mathrm H^\bullet T(L^\bullet)=\mathrm R^\bullet T(A)$.
\end{lemma}

\begin{proof}
C'est un cas particulier de (T, 2.5.2).
\end{proof}

\begin{env}[11.5.2.5]
\label{0.11.5.2.5-fr}
La démonstration de (\hyperref[0.11.5.2-fr]{11.5.2}) se termine maintenant de
façon immédiate, car $L'^{\bullet\bullet}$ est une résolution injective de
$K^\bullet$ dans la catégorie abélienne $\mathcal K$, autrement dit,
$K^\bullet\mapsto\mathrm H^\bullet(T(L'^{\bullet\bullet}))$ n'est autre que
le foncteur cohomologique dérivé droit de $T$ dans la catégorie $\mathcal K$
(T, 2.3).
\end{env}

\begin{proposition}[11.5.3]
\label{0.11.5.3-fr}
Sous les hypothèses de (\hyperref[0.11.5.1-fr]{11.5.1}) concernant
$\mathcal C$, $\mathcal C'$ et $T$, soit
$L^{\bullet\bullet}=(L^{ij})$ un bicomplexe tel que $L^{ij}=0$ pour $j<0$ et
que, pour tout $i$, $L^{i,\bullet}$ soit une résolution de $K^i$~; supposons
enfin que $\mathrm R^nT(L^{ij})=0$ pour tout couple $(i,j)$ et tout $n>0$.
Alors il existe un isomorphisme fonctoriel
\[
  \label{0.11.5.3.1-fr}
  \mathrm R^\bullet T(K^\bullet)\xrightarrow{\sim}
  \mathrm H^\bullet(T(L^{\bullet\bullet})).
  \tag{11.5.3.1}
\]
\end{proposition}

\begin{proof}
Soit $L_{(r)}^{\bullet\bullet}=(L_{(r)}^{ij})$ le bicomplexe tel que
$L_{(r)}^{ij}=0$ pour $i<r$, $L_{(r)}^{ij}=L^{ij}$ pour $i\geq r$~; il est
immédiat que $L^{\bullet\bullet}$ est limite inductive de
$L_{(r)}^{\bullet\bullet}$ pour $r$ tendant vers $-\infty$~; en vertu de
l'hypothèse sur $T$ et de (\hyperref[0.11.5.1-fr]{11.5.1}), il suffit donc de
prouver la proposition lorsque $K^\bullet$ est limité inférieurement, par
exemple $K^i=0$ pour $i<0$, et $L^{ij}=0$ pour $i<0$. Soit alors
$L'^{\bullet\bullet}=(L'^{ij})$ une résolution droite de $K^\bullet$ formée
d'objets de $\mathfrak J$ (\hyperref[0.11.5.2.2-fr]{11.5.2.2})~; il y a un
morphisme de bicomplexes
$L^{\bullet\bullet}\to L'^{\bullet\bullet}$ compatible avec les
augmentations, d'où un morphisme
$'E(T(L^{\bullet\bullet}))\to{}'E(T(L'^{\bullet\bullet}))$ pour les premières
suites spectrales~; le lemme (\hyperref[0.11.5.2.4-fr]{11.5.2.4}) montre,
comme dans (\hyperref[0.11.5.2.3-fr]{11.5.2.3}), que ce morphisme est un
\emph{isomorphisme}, d'où la conclusion.
\end{proof}

\begin{remark}[11.5.4]
\label{0.11.5.4-fr}
Les raisonnements précédents prouvent que les conclusions de
(\hyperref[0.11.5.2-fr]{11.5.2}) et
(\hyperref[0.11.5.3-fr]{11.5.3}) sont valables dans la catégorie des
complexes $K^\bullet$ limités inférieurement, \emph{sans supposer que $T$
permute aux limites inductives filtrantes}. En outre, quand on considère
seulement la catégorie $\mathcal K$ des complexes $K^\bullet$ tels que
$K^i=0$ pour $i<0$, la caractérisation de $\mathrm R^\bullet T(K^\bullet)$
comme le système des \emph{foncteurs dérivés droits} de $T$ dans $\mathcal K$
montre que ce foncteur cohomologique est \emph{universel} (T, 2.3).

Un autre cas où (\hyperref[0.11.5.2-fr]{11.5.2}) est valable sans faire
d'hypothèse supplémentaire sur $T$ est le cas où il existe un entier $m>0$
tel que tout objet de $\mathcal C$ admette une résolution injective de longueur
$\leq m$. En effet, dans la démonstration de
(\hyperref[0.11.5.1-fr]{11.5.1}), toutes les résolutions injectives d'objets
de $\mathcal C$ qui interviennent peuvent être prises de longueur $\leq m$,
d'où il résulte aussitôt que les termes de degré total $n$ du bicomplexe
$T(L_{(r)}^{\bullet\bullet})$ sont égaux à ceux de
$T(L^{\bullet\bullet})$ et en nombre fini, dès que $r$ est assez grand, ce qui
entraîne que pour tout $n$,
$\mathrm H^n(T(L^{\bullet\bullet}))=
\mathrm H^n(T(L_{(r)}^{\bullet\bullet}))$ dès que $r$ est assez grand. Avec
les notations de (\hyperref[0.11.5.2-fr]{11.5.2}), on a donc aussi
$\mathrm R^nT(K_{(r)}^\bullet)=\mathrm R^nT(K^\bullet)$ pour $r$ assez grand
(dépendant de $n$) et de même pour $K'^\bullet$ et $K''^\bullet$, d'où la
conclusion. De la même manière, (\hyperref[0.11.5.3-fr]{11.5.3}) est valable
sans condition supplémentaire sur $T$ lorsque $\mathcal C$ vérifie
l'hypothèse précédente et que l'on suppose que les résolutions
$L^{i,\bullet}$ sont de longueur $\leq m$.
\end{remark}

\begin{env}[11.5.5]
\label{0.11.5.5-fr}
Le résultat de (\hyperref[0.11.5.2-fr]{11.5.2}) se généralise aux
multifoncteurs covariants. Considérons par exemple la situation de
(\hyperref[0.11.4.6-fr]{11.4.6}), où l'on suppose que dans $\mathcal C$,
$\mathcal C'$ et $\mathcal C''$
